\vspace{2.5pt plus 3pt minus 0.5pt}\noindent\hspace*{1.2em}\kbvnum{6.}\hspace{0.55em}ಕೃದಂತ ನಾಮಪದಗಳೆಂದರೇನು? ಉದಾಹರಿಸಿ.\par

\vspace{2.5pt plus 3pt minus 0.5pt}\noindent\hspace*{1.2em}\kbvnum{7.}\hspace{0.55em}ಕೃದಂತ ನಾಮಗಳು ನಾಮಪದಗಳನ್ನು ವಿಶೇಷಿಸುತ್ತಿರುವುದು ಎಂಬುದನ್ನು ಸಮರ್ಥಿಸಲು ಎರಡು ಉದಾಹರಣೆಗಳನ್ನು ಕೊಡಿ.\par

\vspace{2.5pt plus 3pt minus 0.5pt}\noindent\hspace*{1.2em}\kbvnum{8.}\hspace{0.55em}ಕೃದಂತ ಭಾವನಾಮ ಎಂದರೇನು? ಉದಾಹರಣೆ ಸಹಿತ ತಿಳಿಸಿ.\par

\vspace{2.5pt plus 3pt minus 0.5pt}\noindent\hspace*{1.2em}\kbvnum{9.}\hspace{0.55em}ಮಾಡು, ಕೊರೆ, ಓಡು, ನಂಬು, ಅಳೆ—ಈ ಧಾತುಗಳಿಂದ ಹುಟ್ಟುವ ಕೃದಂತ ಭಾವನಾಮಗಳನ್ನು ತಿಳಿಸಿರಿ.\par

\vspace{2.5pt plus 3pt minus 0.5pt}\noindent\hspace*{1.2em}\kbvnum{10.}\hspace{0.55em}“ಹೊಗಳುವ” ಈ ಕೃದಂತ ನಾಮಪ್ರಕೃತಿಗೆ ವಿಭಕ್ತಿ ಪ್ರತ್ಯಯಗಳನ್ನು ಸೇರಿಸಿ ಬರೆಯಿರಿ.\par

\vspace{2.5pt plus 3pt minus 0.5pt}\noindent\hspace*{1.2em}\kbvnum{11.}\hspace{0.55em}ಇರಿತ, ಉಡಿಗೆ—ಈ ಶಬ್ದಗಳನ್ನು ನಿಮ್ಮ ಸ್ವಂತ ವಾಕ್ಯಗಳಲ್ಲಿ ಬಳಸಿರಿ.\par

\vspace{2.5pt plus 3pt minus 0.5pt}\noindent\hspace*{1.2em}\kbvnum{12.}\hspace{0.55em}ಕೃದಂತಾವ್ಯಯ ಎಂದರೇನು? ಉದಾಹರಣೆಗಳ ಸಹಿತ ತಿಳಿಸಿ.\par

\vspace{2.5pt plus 3pt minus 0.5pt}\noindent\hspace*{1.2em}\kbvnum{13.}\hspace{0.55em}ಕೃದಂತಾವ್ಯಯಗಳ ಲಕ್ಷಣಗಳೇನು?\par

\vspace{2.5pt plus 3pt minus 0.5pt}\noindent\hspace*{1.2em}\kbvnum{14.}\hspace{0.55em}“ಬೀದಿಯಲ್ಲಿ ಬರುತ್ತ ಮುಗ್ಗರಿಸಿ ಬಿದ್ದನು.” ಈ ವಾಕ್ಯದಲ್ಲಿ ಕೃದಂತಾವ್ಯಯ ಯಾವುದು?\par

\vspace{2.5pt plus 3pt minus 0.5pt}\noindent\hspace*{1.2em}\kbvnum{15.}\hspace{0.55em}“ಮಾಡಲಿಕ್ಕೆ” ಈ ಶಬ್ದದಲ್ಲಿ ಧಾತು ಯಾವುದು? ಮತ್ತು ಪ್ರತ್ಯಯ ಯಾವುದು?\par

\vspace{2.5pt plus 3pt minus 0.5pt}\noindent\hspace*{1.2em}\kbvnum{16.}\hspace{0.55em}ಓಟ, ಹೇಸಿಗೆ—ಈ ಪದಗಳಲ್ಲಿರುವ ಮೂಲ ಧಾತುಗಳನ್ನು ಗುರ್ತಿಸಿ.\par

\vspace{2.5pt plus 3pt minus 0.5pt}\noindent\hspace*{1.2em}\kbvnum{17.}\hspace{0.55em}ಭೂತ ಕೃದಂತ, ಭವಿಷ್ಯತ್‌ ಕೃದಂತ, ನಿಷೇಧ ಕೃದಂತ—ಇವು ಗಳಿಗೆ ಒಂದೊಂದು ಉದಾಹರಣೆಯನ್ನು ಕೊಡಿ.\par

\kbvchapter{ಪಾಠ ೧೫}{ತದ್ಧಿತಾಂತಗಳು}

\indent ಆಟವನ್ನು ಆಡುವವನು ಆಟಗಾರ; ಇಲ್ಲಿ “ಆಟ” ಎಂಬ ನಾಮ ಪ್ರಕೃತಿಯೊಂದಿಗೆ “ಗಾರ” ಎಂಬ ಪ್ರತ್ಯಯ ಸೇರಿ “ಆಟಗಾರ” ಆಗಿರುವುದು. “ಗಾರ” ಪ್ರತ್ಯಯವನ್ನು “ಆಡುವವನು” ಎಂಬ ಅರ್ಥದಲ್ಲಿ ಬಳಸಲಾಗಿದೆ.\par

\indent ಆಟವನ್ನು + ಗಾರ = ಆಟ + ಅನ್ನು + ಗಾರ; ಆದರೆ, ಒಂದೇ ನಾಮ ಪ್ರಕೃತಿಯ ಮೇಲೆ ಎರಡು ಪ್ರತ್ಯಯಗಳು ಬರುವುದು ಅಸಂಭವ. ಈ ಕಾರಣ\kbvcontinues
